Write a program that calculates the \textbf{sum of the integers from} \(\boldsymbol{0}\) \textbf{to} \(\boldsymbol{n}\) in two ways, with a loop and with a mathematical formula, and checks that both agree.

\subsection{Environment setup}

\begin{enumerate}

    \item \textbf{Install and configure the IDE}: download and install it from \url{https://www.jetbrains.com/idea/download/}.

    \item \textbf{Create the project structure}:

        \begin{enumerate}

            \item Select \uibutton{New Project} from the welcome screen.

            \item Choose \uibutton{Java} in the left menu.

            \item Configure the project details:

                \begin{itemize}

                    \item \textbf{Project name}: \uibutton{Laboratory assignment 0}.

                    \item \textbf{Project location}: choose a folder where you can save and later delete your work.

                \end{itemize}

            \item Click the \uibutton{Create} button to generate the workspace.

        \end{enumerate}

\end{enumerate}

\subsection{Implementation of the summation methods}

Once the project workspace has been generated, remove the \texttt{Main.java} file inside the source (\texttt{src}) folder of your project and copy the \texttt{Sum.java} template provided with this assignment into it. Each \texttt{TODO} comment in the template marks a part you have to write.

In Java, functions are called \textbf{methods}. You declare one as a function in C, with the \texttt{static} keyword in front, for example \texttt{static int sum1(int n)}. Java requires it here, as \texttt{main} calls these methods directly; for now, write it as it is: what it means will be covered in later sessions.

\begin{enumerate}

    \item \textbf{Iterative sum}: write a method named \texttt{sum1} that receives an \texttt{int} \(n\) and returns the sum of the integers from \(0\) to \(n\): \(sum1(n) = \sum_{i = 0}^{n} i\).

        \begin{itemize}

            \item \textbf{Logic}: translate the implementation from C. Initialize a variable to \(0\) and use a \texttt{for} loop to add the values from \(0\) to \(n\).

        \end{itemize}

    \item \textbf{Formula sum}: write a second method named \texttt{sum2} that returns the same sum using the arithmetic progression formula: \(sum2(n) = \frac{n(n + 1)}{2}\).

        \begin{itemize}

            \item \textbf{Logic}: its body is a single \texttt{return} statement.

        \end{itemize}

    \item \textbf{Verification}: write the body of \texttt{main} to test both methods, printing the results with \texttt{System.out.printf}, which works as \texttt{printf} in C (with \texttt{\%d} for an \texttt{int} and \texttt{\textbackslash{}n} to end the line).

        \begin{enumerate}

            \item \textbf{Individual test}: print the results of \texttt{sum1} and \texttt{sum2} for \(n = 11\). Both must be \(66\).

            \item \textbf{Batch test}: use a loop to go through the following array of values: \texttt{\{3, 4, 13, 21, 67, 102, 155, 365, 1007\}}. For each value, print \(n\), the results of both methods and whether they are equal.

            \item \textbf{Automatic comparison}: instead of checking the numbers yourself, let the program tell you. Compare the results of \texttt{sum1} and \texttt{sum2} with \texttt{==}, which gives a \texttt{boolean}, and print it with \texttt{\%b}, which prints \texttt{true} or \texttt{false}. Every value in the array must print \texttt{true}.

        \end{enumerate}

    \item \textbf{Documentation comments}: document the \texttt{Sum} class and the \texttt{sum1} and \texttt{sum2} methods with \textbf{Javadoc comments} (\texttt{/** ... */}), placed right above what they describe. The comment of \texttt{main} in the template is an example to follow:

        \begin{itemize}

            \item The first sentence says what the class or method does.

            \item \texttt{@param} followed by a parameter's name describes that parameter, one line per parameter.

            \item \texttt{@return} describes the value the method returns.

            \item \texttt{@author} followed by your name, in the comment of the class.

            \item Javadoc comments are HTML, not Markdown: write code (names, expressions) as \texttt{\{@code sum1\}} and use HTML for formulas, such as \texttt{<i>n</i>(<i>n</i> + 1) / 2}.

        \end{itemize}

    \item \textbf{Documentation generation}: generate the documentation of your project to see what your comments become.

        \begin{enumerate}

            \item Open the \uibutton{Tools} menu.

            \item Select \uibutton{Generate JavaDoc...}.

            \item Set the output directory to your project folder followed by \texttt{/doc}.

            \item Once the process is complete, check that a new folder named \texttt{doc} has been created in your project directory.

            \item Open its \texttt{index.html} file in a web browser. Observe how your Javadoc comments have been turned into an HTML documentation page, listing the class, its methods, their parameters and their return values.

        \end{enumerate}

\end{enumerate}
